\documentclass[11pt]{amsart}
\usepackage[T1]{fontenc}
\usepackage{lmodern}
\usepackage[margin=1in]{geometry}
\usepackage{amsmath,amssymb,amsthm,booktabs,array,xcolor}
\usepackage{microtype}
\usepackage[colorlinks=true,linkcolor=blue!45!black,citecolor=blue!45!black,urlcolor=blue!45!black]{hyperref}
\newtheorem{theorem}{Theorem}[section]
\newtheorem{lemma}[theorem]{Lemma}
\newtheorem{proposition}[theorem]{Proposition}
\newtheorem{corollary}[theorem]{Corollary}
\theoremstyle{definition}
\newtheorem{definition}[theorem]{Definition}
\theoremstyle{remark}
\newtheorem{remark}[theorem]{Remark}
\newcommand{\bs}{\operatorname{bs}}
\newcommand{\wt}{\operatorname{wt}}
\newcommand{\dist}{\operatorname{dist}}
\newcommand{\ind}{\mathbf1}
\newcommand{\old}{\mathrm{old}}
\newcommand{\new}{\mathrm{new}}
\newcommand{\ceil}[1]{\left\lceil#1\right\rceil}
\hypersetup{pdftitle={First-difference labels and a coupled-profile refinement for a sensitivity construction}}
\title[Explicit labels and refined sensitivity profiles]{First-difference labels and a coupled-profile refinement\texorpdfstring{\\}{} for a sensitivity construction}
\date{15 October 2026}
\begin{document}
\begin{abstract}
The OpenAI family-132 construction of a superquadratic separation between
ordinary sensitivity and block sensitivity uses a labelled regular tournament
with no head-coherent set of sixteen vertices. We give an elementary directed
adaptation of classical first-difference colouring that supplies the exact
labels on an infinite parameter family, with polynomial logarithmic bit time
per edge query and no witness advice. In the indexed cyclic construction,
the coherent-set bound is fourteen throughout the family and exactly seven
for the fixed seed below. These are direct consequences of the same proof.
We also refine the construction's
coupled sensitivity estimates: candidate target repairs usually cross two
child thresholds, and target and gate candidate counts obey joint constraints.
We prove the resulting six-branch ordinary and five-branch joint zero-side
bounds. On one fixed explicit seed, exact arithmetic certifies an improvement
of the amplified exponent bound from
$2.015411246837677710\ldots$ to $2.015411547511314020\ldots$,
using the exact cap seven in both analyses. The former already follows from
immediate refinements of the published proof with these explicit labels.
Independent unrefereed manuscripts posted on 7 and 9 October 2026 report
larger exponents, about $2.0321$ and $2.0271$, using different parameter
schedules; the second also proves part of the target/gate counting used
below.
The contribution is a deterministic instantiation and a local proof
refinement, not a new separation architecture or a record-exponent claim.
\end{abstract}
\maketitle

\section{Scope and attribution}
For a total Boolean function $f:\{0,1\}^n\to\{0,1\}$, write $x^D$ for
the input obtained by flipping the coordinates in $D$. Its \emph{sensitivity}
$s(f,x)$ counts coordinates $v$ with $f(x^{\{v\}})\ne f(x)$.
Its \emph{block sensitivity} $\bs(f,x)$ is the largest number of pairwise
disjoint nonempty sets $D$ whose separate flips change $f(x)$. Put
$s(f)=\max_xs(f,x)$ and $\bs(f)=\max_x\bs(f,x)$.

The pinned OpenAI manuscript~\cite{OpenAI132} proves that
$\bs(f)/s(f)^2$ is unbounded for total Boolean functions and obtains a
fixed exponent greater than two by composition. Its construction adapts
gated tournaments from Meiburg~\cite{Meiburg}; the latter's headline
comparison is with spectral sensitivity $\lambda$, not ordinary
sensitivity $s$. Since $\lambda\le s$, a lower bound in terms of
$\lambda$ does not by itself imply the ordinary-sensitivity separation.
OR balancing and composition are established tools; see
Ambainis--Sun~\cite{AmbainisSun} and Tal~\cite{Tal}.

This note proves two statements within the family-132 framework.
First, a classical first-difference rule has a useful arbitrary-orientation
head-coherence bound, yielding a succinct deterministic replacement for the
source's probabilistic labelling. Second, a finer analysis of the source's
candidate clauses gives stronger coupled profile inequalities. All Boolean
predicates and the zero-block construction remain those of~\cite{OpenAI132}.

First-difference methods, including position-plus-own-bit local asymmetry,
are classical: see Cole--Vishkin~\cite{ColeVishkin}, Vishkin's
notes~\cite{VishkinNotes}, and the hypercube colouring in
Gishboliner--Milojevi\'c--Sudakov--Wigderson~\cite{GMSW}.
The directed application below is an elementary adaptation, and historical
novelty of its exact formulation is not assured. The numerical comparison
later in the note is deliberately against an immediate refinement of the
published proof on the \emph{same seed}. An exponent exceeding $2.0154$
already follows from that baseline.
The sharper cyclic coherence bound and its exact seven-vertex witness are
elementary corollaries; they are not separate methodological novelty claims.

\emph{Related refinements.} Two AI-assisted, unrefereed manuscripts in the
ProveIt repository refine the same source construction~\cite{ProveItI,ProveItII}.
We have not independently verified them, and summarize only their stated
results. The first, dated 7 October 2026, allows level-dependent tournament
sizes with four positions per row and reports the exponent
$2.0320827244\ldots$ for ordinary sensitivity. Its zero-side recurrence already contains the
gate envelope $\max\{b+e,3e\}$ that we include in the baseline of
Section~\ref{sec:baseline}. It obtains labels from a union bound and lists
explicit labellings as an open problem. The second, a draft dated 8 October
and posted 9 October, reports the exponent $3667/1809\approx2.0270868$ and a
deterministic conditional-expectation labelling that is impractical at its
parameters. Its joint target/gate budget lemma shows that two gate
candidates leave at most $t-2$ target candidates, and three leave at most
$t-4$. With $t=16$ these are the target counts in our fifth and sixth cases
in Section~\ref{sec:cases}. Relative to those manuscripts, our additional profile
ingredients are the charging of all but at most one target repair to a
joint profile, the coherent unions in the one-gate and ordinary-sink cases,
and the stronger-index joint candidate lists. Our labels serve the
family-132 regime with about $\sqrt M$ labels; they do not directly supply
the four-label exclusion condition of~\cite{ProveItI}. The fixed-seed
comparison in Section~\ref{sec:numerical} is not competitive with the
exponents above and is not offered as one.

\section{Explicit labels for the exact tournament condition}
\begin{definition}[Head coherence]\label{def:coherence}
A tournament directs exactly one edge between each pair of distinct vertices.
Given labels $\ell(i,j)$ on its edges $i\to j$, a vertex set $I$ is
\emph{head-coherent} if there are labels $m_j$, $j\in I$, such that
\[
 \ell(i,j)=m_j\qquad(i,j\in I,\ i\to j).
\]
A vertex with no incoming edge in $I$ imposes no restriction on $m_j$.
\end{definition}

The condition required in~\cite[Lemma 2.1]{OpenAI132} is a regular
tournament on
\begin{equation}\label{eq:parameters}
 k=2M^2+1,\qquad r=\ceil{\sqrt M},\qquad M\ge3,
\end{equation}
with outdegree $M^2$, labels in $\{0,\ldots,r-1\}$, and no head-coherent
sixteen-set. Independent labels give the source's existence bound
$k^{16}r^{-104}\le3^{16}M^{-20}<1$.

\begin{theorem}[First-difference bound]\label{thm:digits}
Let $h\ge1$ and $q\ge2$. Give any subset of $\{0,\ldots,q-1\}^h$ any
tournament orientation. For $x\to y$, let $u$ be the first position at
which $x_u\ne y_u$, and set
\begin{equation}\label{eq:label}
 \ell(x,y)=(u-1)(q-1)+x_u-\ind\{x_u>y_u\}.
\end{equation}
These $h(q-1)$ labels have no head-coherent set of more than $h+2$ vertices.
\end{theorem}
\begin{proof}
At a fixed head digit $b$, the map $a\mapsto a-\ind\{a>b\}$ is a
bijection from the digits other than $b$ to $\{0,\ldots,q-2\}$.
Thus equal labels at a fixed head mean equal first-difference positions
and equal source digits there. Positions have disjoint palettes.

Partition a coherent set by first digit into buckets. A vertex receiving
an edge from its own bucket has an incoming label with position at least
two. It cannot also receive an edge from outside that bucket, whose
position would be one. It therefore sends an edge to every outside vertex.
Every bucket of size at least two contains such a vertex. Two buckets of
size at least two would consequently require both directions between a
chosen vertex of each. Hence there is at most one large bucket $B$.

If $B$ exists, choose $z\in B$ receiving an edge from $B$. Every vertex
outside $B$ receives an edge from $z$, and all outside buckets are singletons.
If two outside vertices existed, their mutual edge would give its head
two position-one labels with distinct source digits, one from $B$ and one
from the other singleton. Fixed-head injectivity excludes this. There is
therefore at most one vertex outside $B$.

If all buckets are singletons, each vertex has indegree at most one;
summing indegrees gives $|I|(|I|-1)/2\le|I|$, so $|I|\le3$.
Inside a large bucket, delete its common first digit and subtract the
palette offset. Coherence is preserved. The one-digit bound is three,
and the preceding argument gives $C_h\le\max\{3,C_{h-1}+1\}$ for the
maximum possible size. Induction yields $C_h\le h+2$.
\end{proof}

\begin{remark}\label{rem:sharp}
For $q\ge3$ the bound is sharp over arbitrary orientations. Start with
three words differing only in the last digit, oriented as a directed
triangle. Successively add a word with a nonzero digit one position
earlier than all preceding nonzero digits, directing all preceding
vertices toward it. Each new head receives one label, giving $h+2$
coherent vertices. The proof also works for unequal radices $q_u$,
using $\sum_u(q_u-1)$ labels. If the last radix is two, the base bound
is two and the resulting bound is $h+1$.
\end{remark}

\begin{theorem}[Exact infinite family]\label{thm:generator}
For every integer $q\ge3$, let
\begin{equation}\label{eq:family}
 M=[13(q-1)]^2,\qquad r=13(q-1),\qquad k=2M^2+1.
\end{equation}
There is a uniform deterministic tournament on $k$ vertices, of outdegree
$M^2$, with labels in $\{0,\ldots,r-1\}$ and no head-coherent sixteen-set.
Each edge orientation and label
can be evaluated using $O(\log M)$ storage bits and
$O((\log M)^2)$ bit operations, with no precomputed witness or edge table.
\end{theorem}
\begin{proof}
On $\{0,\ldots,k-1\}$ use the cyclic orientation
\begin{equation}\label{eq:cyclic}
 i\to j\quad\Longleftrightarrow\quad
 1\le(j-i)\bmod k\le M^2.
\end{equation}
There are exactly $M^2$ successors per vertex, and odd $k$ ensures one
direction per pair. Encode each index with thirteen base-$q$ digits,
including leading zeros, and use~\eqref{eq:label}. This uses
$13(q-1)=r=\ceil{\sqrt M}$ labels.

There are enough words: for $q=3$,
$k=913953<1594323=3^{13}$. For $q\ge4$,
\[
 q^{13}-k>q^4(q^9-2\cdot13^4)-1>0,
\]
since $4^9=262144>57122=2\cdot13^4$.
Theorem~\ref{thm:digits} bounds coherent sets by fifteen.
All integers involved have $O(\log M)$ bits. A fixed number of ordinary
integer operations extracts thirteen digits and evaluates
\eqref{eq:label}--\eqref{eq:cyclic}; schoolbook arithmetic gives the bit
bound. The algorithm is uniform in $q$, and the family is unbounded.
\end{proof}

\begin{corollary}[Cyclic coherence and the exact seven-bound]
\label{cor:cyclic}
For the indexed tournament in Theorem~\ref{thm:generator}, put
$H=\ceil{\log_q k}\le13$. Every head-coherent set has at most $H+1$
vertices. The uniform bound fourteen is attained at $q=3$.
At $q=1001$, the maximum coherent-set size is exactly seven.
\end{corollary}
\begin{proof}
Delete the $13-H$ common leading zero digits and subtract the resulting
common palette offset. Within a fixed prefix of length $H-1$, indices span
at most $q-1\le13^4(q-1)^4=M^2$. Thus every edge in such a final-digit
bucket points from the smaller index to the larger. Three distinct digits
would give the largest vertex two incoming edges with different source
digits, contradicting fixed-head injectivity. The one-digit coherent bound
is therefore two.

The bucket argument of Theorem~\ref{thm:digits} applies inside every prefix.
There is at most one nonsingleton bucket, with at most one vertex outside
it; if all buckets are singletons there are at most three vertices.
Hence the bounds for suffix length $s$ satisfy
\[
 C_1\le2,\qquad
 C_s\le\max\{3,C_{s-1}+1\}\le s+1\quad(s\ge2).
\]
This gives the bound $H+1$, including its separate one-digit base case.

For $q=3$, the fourteen vertices
\[
 \{3^{11}\}\cup\{3^{11}+3^j:0\le j\le10\}
 \cup\{2\cdot3^{11},3^{12}\}
\]
are valid indices in $[177147,531441]$. Their span is
$354294<M^2=456976$, so all edges point increasingly. At each new head,
the earlier vertices have a common first-difference position and source
digit: first the successive suffix powers, then the digit one against
two at place $3^{11}$, then zero against one at place $3^{12}$.
This is a coherent fourteen-set.

For $q=1001$, exact arithmetic gives
\[
 q^5=1{,}005{,}010{,}010{,}005{,}001
 <M^2=28{,}561{,}000{,}000{,}000{,}000
\]
and $q^5<k<q^6$, so $H=6$ and the upper bound is seven.
The seven valid vertices
\[
 I=\{0,1,q,q^2,q^3,q^4,q^5\}
\]
span less than $M^2$, so every edge points increasingly. For each head
$q^j$, every earlier vertex agrees before its leading one and has zero
in that position. All incoming labels at that head coincide. This proves
the matching lower bound seven.
\end{proof}

\begin{remark}[Extent of the classical correspondence]
Uncompressed labels $(u,x_u)$ use the classical position-plus-source-digit
idea. Compression in~\eqref{eq:label} preserves equality at each
\emph{fixed head}, which is precisely what Definition~\ref{def:coherence}
uses. It need not preserve a global proper-colouring or walk-colouring
property. The bucket argument proves the directed bound separately.
\end{remark}

\section{The inherited Boolean recursion}
Fix $d\ge1$, $M\ge3$, and a labelled regular tournament with the parameters
in~\eqref{eq:parameters} and no head-coherent sixteen-set.
Here $d$ is recursion depth; the digit length $h=13$ is independent of it.
Put
\begin{equation}\label{eq:heights}
 L=d+2,\qquad H_j=d+1-j,\qquad N_j=LM(kr)^j
 \quad(0\le j\le d).
\end{equation}
The following predicates are exactly those of~\cite{OpenAI132}.
Partition the base coordinates into $M$ disjoint $L$-sets $E_v$, let
$z_v=\ind_{E_v}$, and set
\begin{equation}\label{eq:base}
 D(x)=\min_v\dist(x,z_v),\qquad
 P_{0,p}(x)=\ind\{D(x)\le d+1-p\},\quad 1\le p\le H_0.
\end{equation}
At level $j\ge1$, arrange the input into $kr$ children $x_{i,t}$,
where $0\le i<k$ and $0\le t<r$. Write $Q=H_{j-1}$. For
$1\le p\le H_j$, clause $i$ for predicate index $p$ requires
\begin{equation}\label{eq:clauses}
 \begin{aligned}
 P_{j-1,Q}(x_{i,t})&=1 &&(0\le t<r),\\
 P_{j-1,Q-p}(x_{v,\ell(i,v)})&=0 &&(i\to v).
 \end{aligned}
\end{equation}
Let $P_{j,p}$ be the OR of all these clauses. The first line consists
of \emph{targets}, and the second of \emph{gates}. Each condition in a
fixed clause uses a distinct child.

\begin{lemma}[Basic properties and zero blocks; inherited]\label{lem:blocks}
These predicates are total and nested:
$P_{j,1}\ge\cdots\ge P_{j,H_j}$. All reject zero. There are $Mk^j$
pairwise disjoint nonempty blocks of size $Lr^j$ whose separate flips
at zero make every level-$j$ predicate accept.
\end{lemma}
\begin{proof}
All indices in~\eqref{eq:clauses} are valid, since $1\le Q-p<Q$.
The formulas therefore define a value on every input. Nesting holds at
the base, and increasing $p$ decreases the gate index, strengthening the
requirement that a nested predicate be zero. Targets are unchanged.

At the base, $D(0)=L>d$, and flipping $E_v$ produces a centre, where
all predicates accept. Inductively every target at zero is false.
Given a preceding sensitive block, copy it into all $r$ children of
one row and take the union. Those targets accept after the flip, while
every other row remains zero, so every outgoing gate passes. Distinct
rows or preceding block indices give disjoint unions. Their number is
$k\cdot Mk^{j-1}$ and their size is $r\cdot Lr^{j-1}$.
\end{proof}

\begin{definition}[Ordinary and joint profiles]\label{def:profiles}
For $b\in\{0,1\}$, let
\[
 S_{j,b}=\max_{p,x:\,P_{j,p}(x)=b}
 \bigl|\{v:P_{j,p}(x^{\{v\}})=1-b\}\bigr|.
\]
The joint profile is
\[
 J_{j,b}=\max_{\substack{p<p',\,x:\\P_{j,p}(x)=P_{j,p'}(x)=b}}
 \bigl|\{v:P_{j,p}(x^{\{v\}})=P_{j,p'}(x^{\{v\}})=1-b\}\bigr|.
\]
Both predicates use the same input and the same raw coordinate. Empty
maxima are zero; in particular $J_{d,0}=J_{d,1}=0$. Always
$0\le J_{j,b}\le S_{j,b}$.
\end{definition}

\begin{lemma}[Base profiles; inherited]\label{lem:baseprofiles}
At level zero,
\[
 S_{0,0}\le d+1,\qquad S_{0,1}\le LM,\qquad J_{0,0}=J_{0,1}=0.
\]
The displayed source recurrence uses the looser bound $S_{0,0}\le L$.
\end{lemma}
\begin{proof}
Distinct centres have distance $2L=2d+4$, so at most one centre can
lie within distance $d+1$ of a rejecting input. A successful flip must
move toward that centre from distance $R+1$, where $R\le d$ is the
predicate's radius. There are at most $R+1\le d+1$ such coordinates.
The one-side bound is the input dimension. Finally $D$ is 1-Lipschitz
on adjacent inputs, so one flip cannot cross two distinct integer radii.
\end{proof}

\section{A coupled candidate refinement}
At a fixed level write the preceding actual profiles as
\begin{equation}\label{eq:childnames}
 a=S_{j-1,0},\qquad b=S_{j-1,1},\qquad
 c=J_{j-1,0},\qquad e=J_{j-1,1}.
\end{equation}
Thus $c\le a$ and $e\le b$. These letters are separate from the final
seed sensitivity bound $A$ and zero-block bound $B$.

\begin{theorem}[Refined coupled profiles]\label{thm:profiles}
For the predicates above and any admissible labels,
\begin{align}
 S_{j,0}&\le F_0(a,b,c,e),& J_{j,0}&\le F_J(a,b,c,e),\label{eq:newzero}\\
 S_{j,1}&\le M^2a+rb,& J_{j,1}&\le M^2c+rb,\label{eq:ones}
\intertext{where}
 F_0(a,b,c,e)&=\max\left\{
 \begin{array}{lll}
 a+14c,&a+13c+b,&15c+e,\\
 13c+b+e,&14c+2e,&12c+3e
 \end{array}\right\},\label{eq:F0}\\
 F_J(a,b,c,e)&=\max\left\{
 \begin{array}{lll}
 a+14c,&a+13c+e,&15c+e,\\
 14c+2e,&12c+3e
 \end{array}\right\}.\label{eq:FJ}
\end{align}
All branches have nonnegative coefficients. Consequently these bounds
can be propagated from componentwise upper bounds, with optional
definitional capping $J_{j,b}\le S_{j,b}$.
\end{theorem}

The one-side inequalities~\eqref{eq:ones} are inherited from the source.
The refinement concerns its treatment of candidate clauses at a rejecting
input. We prove every branch, including the joint transition statement.

\subsection{Candidate structure}
Fix a parent index and a rejecting input. Suppress the child-level prefix,
write $Q$ for the strongest child index, and $g<Q$ for the gate index.
Let $T$ be the clauses with exactly one failed target and no failed gate.
Let $G$ be the clauses with exactly one failed gate and no failed target.
Every successful coordinate repairs a clause in these fixed lists: one
raw bit can affect at most one condition of a clause. Write $m_t$ for
the missing-target position of a row $t\in T$.

\begin{lemma}\label{lem:candidates}
The set $T$ is head-coherent and $|T|\le15$. Every edge between $G$ and
$T$ points from $G$ to $T$. Moreover $|G|\le3$; if $|G|=3$, it is a
directed cycle and all three failed gates are internal to $G$.
\end{lemma}
\begin{proof}
For $s\to t$ inside $T$, the passing gate gives $P_g=0$, hence
$P_Q=0$, at its selected child in row $t$. The unique such child is
$m_t$, so $\ell(s,t)=m_t$. This proves coherence. An edge from $T$ to
$G$ would fail its gate because all children of a $G$ row have
$P_Q=1$, hence $P_g=1$. Thus no such edge exists.

Every internal outgoing edge of a $G$ vertex is a failed gate, so each
has internal outdegree at most one. Counting edges yields
$|G|(|G|-1)/2\le|G|$, hence $|G|\le3$. With three vertices the only
possibility is a directed cycle, using each vertex's unique failure.
\end{proof}

Call a target candidate \emph{ordinary} if its missing target has
$P_g=1$ and $P_Q=0$, and \emph{joint} otherwise. An ordinary target
has at most $a$ repair coordinates. A joint target starts with
$P_g=P_Q=0$ and a successful repair makes both one, so it has at most
$c$ repair coordinates. Call a gate candidate ordinary if its failed
gate child has $P_Q=0$, and joint otherwise. An ordinary gate costs
at most $b$ repairs. A joint gate starts with $P_g=P_Q=1$ and repairing
$P_g$ makes both zero, so it costs at most $e$.

\begin{lemma}\label{lem:ordinary}
There is at most one ordinary target. If $|G|\ge2$, every target is
joint. At most one gate candidate is ordinary, and all gate candidates
are joint when $|G|=3$.
\end{lemma}
\begin{proof}
An ordinary target has no incoming edge from $T$: such an edge would
force $P_g=0$ at its missing target. A tournament has at most one
vertex of indegree zero. Also an ordinary target's entire row has
$P_g=1$. Every $G$ vertex therefore uses its unique failed gate on
that row. When $|G|\ge2$, its internal edge already uses one such
failure, a contradiction. Finally every failed gate internal to $G$
is joint, because all targets of its head row pass. With two gate
candidates only the sink could be ordinary; with three, every failure
is internal by Lemma~\ref{lem:candidates}.
\end{proof}

\subsection{All six ordinary zero cases}\label{sec:cases}
The following table gives upper bounds for the combined repair cost.
The possible overlap of repair coordinates only makes the union smaller.
\begin{center}
\small
\begin{tabular}{@{}c p{6.8cm} c@{}}
\toprule
$|G|$ & Condition & Repair bound\\
\midrule
0 & No gate candidate & $a+14c$\\
1 & An ordinary target or an ordinary gate exists & $a+13c+b$\\
1 & Neither ordinary type exists & $15c+e$\\
2 & The sink's failed gate is ordinary & $13c+b+e$\\
2 & Both failed gates are joint & $14c+2e$\\
3 & The gate candidates form a directed cycle & $12c+3e$\\
\bottomrule
\end{tabular}
\end{center}
For completeness, the cardinality argument for each row follows.

\emph{No gate candidate.}
There are at most fifteen targets and at most one is ordinary. Since
$c\le a$, their cost is bounded by $a+14c$, including when no ordinary
target exists or $T$ is empty.

\emph{One gate candidate, an ordinary type present.}
Suppose first that the failed gate is ordinary. If it points into $T$,
its child has $P_Q=0$ and so is that row's missing target; its label
is $m_t$. Every other edge from $G$ into $T$ passes and also has label
$m_t$. Thus $T\cup G$ is head-coherent, whether or not the failure
points into $T$, and $|T|\le14$.

Suppose instead that an ordinary target exists. Its entire row has
$P_g=1$, so the sole gate candidate's unique failure points into that
row. The ordinary target is a source inside $T$ and has just one
incoming edge in $T\cup G$. Assign its head label to that edge's label;
no internal incoming edge competes. All other heads retain their
coherent labels. Again $T\cup G$ is coherent and $|T|\le14$.
This reassignment of the exceptional head label is necessary; its
missing-target position need not label the failed edge. The target cost
is at most $a+13c$, and the gate cost is at most $b$ since $e\le b$.

\emph{One gate candidate, neither ordinary type present.}
The general bound $|T|\le15$ gives target cost $15c$, and the gate
costs $e$. The union need not be coherent in this case.

\emph{Two gate candidates, ordinary sink gate.}
The tail's failed gate is their internal edge; only the sink can fail
outside $G$. If its ordinary failed gate points into $T$, its child
has $P_Q=0$, so its label is the head's missing position $m_t$.
Every other edge into $T$ passes and has its prescribed label.
The internal edge of $G$ imposes only one incoming label there, and
there are no edges from $T$ to $G$. Hence $T\cup G$ is coherent and
$|T|\le13$. All targets are joint by Lemma~\ref{lem:ordinary}; the
gate cost is at most $b+e$.

\emph{Two gate candidates, both joint.}
Only the sink can have a failure outside $G$. All edges from $G$ into
$T$ therefore have the expected head label except possibly at the head
of that one failure. Delete that target if it exists. The remaining
union is coherent, so $|T|-1+2\le15$ and $|T|\le14$. All target
repairs are joint, giving $14c+2e$.

\emph{Three gate candidates.}
All failures are internal to their directed cycle, whose heads each
have one incoming edge. All edges into $T$ pass. Thus $T\cup G$ is
coherent and $|T|\le12$. All candidates are joint, giving $12c+3e$.

These cases exhaust the possibilities. Taking their maximum and then
maximizing over rejecting inputs proves~\eqref{eq:F0}.

\subsection{All joint zero transitions and the one side}
Fix parent indices $p<p'$ and an input where both predicates are zero.
Their gate indices obey $g'=Q-p'<g=Q-p<Q$.
Form $T$ at the stronger parent index $p'$. Form $G$ there too, but
retain only candidates whose unique failed gate also fails at the
weaker index $p$.

A coordinate making both parents accept satisfies some clause at $p'$.
Before the flip this clause has exactly one failed condition. A target
failure puts it in $T$. If it is a gate failure, every other gate
already passes at $p'$, hence also at $p$. All targets pass as well.
If that same gate did not fail at $p$, the weaker parent would already
accept, a contradiction. Thus this clause lies in the retained $G$.
The fixed lists cover every simultaneous success even when the witness
clause depends on the flipped coordinate.

The previous candidate arguments apply to these lists at gate index
$g'$. Deleting gate candidates changes none of their premises. Target
repairs have the same costs $a,c$. Every retained gate, including a
gate ordinary relative to $P_Q$, now costs at most $e$: its distinct
predicates $P_{g'}$ and $P_g$ start at one, and satisfying the stronger
gate makes both zero. Replace $b$ by $e$ in the six rows. The fourth
row $13c+2e$ is dominated by the fifth, $14c+2e$. The remaining five
rows are exactly~\eqref{eq:FJ}. If there are no distinct parent
indices the joint profile is zero by definition.

For ordinary one-to-zero transitions, fix one initially satisfied
clause. Every successful rejection breaks it, by one of its $r$ targets
or $M^2$ gates. These cost at most $rb$ and $M^2a$, respectively.
For joint rejection at indices $p<p'$, fix a clause satisfied at $p'$;
it also satisfies $p$. Since both parents reject after the flip, this
clause must fail at the weaker index $p$. A target failure costs at
most $b$. A gate failure changes $P_g$ from zero to one; initially
$P_{g'}=0$, and nesting forces $P_{g'}=1$ afterward, so it costs at
most $c$. This proves~\eqref{eq:ones} and completes the proof of
Theorem~\ref{thm:profiles}.

\begin{corollary}[Profiles with an arbitrary coherence cap]
\label{cor:cap-profiles}
Suppose every head-coherent set has size at most $K$, where $K\ge3$.
Then the zero-side bounds in Theorem~\ref{thm:profiles} hold with
\begin{align}
 F_0^{(K)}(a,b,c,e)&=\max\left\{
 \begin{array}{ll}
 a+(K-1)c,&a+(K-2)c+b,\\
 Kc+e,&(K-2)c+b+e,\\
 (K-1)c+2e,&(K-3)c+3e
 \end{array}\right\},\label{eq:FK0}\\
 F_J^{(K)}(a,b,c,e)&=\max\left\{
 \begin{array}{ll}
 a+(K-1)c,&a+(K-2)c+e,\\
 Kc+e,&(K-1)c+2e,\\
 (K-3)c+3e&
 \end{array}\right\}.\label{eq:FKJ}
\end{align}
The one-side bounds~\eqref{eq:ones} are unchanged.
\end{corollary}
\begin{proof}
The six ordinary cases above bound the number of targets respectively by
$K,K-1,K,K-2,K-1,K-3$. In the first two cases at most one target costs
$a$ and the others cost $c$; in the remaining cases all cost $c$.
The exceptional head reassignment in the second case, deletion of at most
one target in the fifth, and coherent union with three gate candidates in
the sixth apply unchanged. These are the only uses of the number fifteen.
The gate costs are as before, giving~\eqref{eq:FK0}. For joint acceptance,
use the stronger-parent lists and retain gates failing at both indices.
Every retained gate costs $e$, so substitute $b=e$ and discard the dominated
fourth branch to obtain~\eqref{eq:FKJ}. All coefficients are nonnegative
because $K\ge3$; hence the same monotone propagation is valid.
\end{proof}

\begin{remark}[What the refinement does and does not assert]
The displayed source bounds are
\begin{equation}\label{eq:published}
 S_{j,0}\le16a+b+3e,\qquad J_{j,0}\le16a+3e,
\end{equation}
with the same one-side bounds. The new argument uses joint target
repairs and coupled counts; it is not just replacement of $16$ by $15$.
For example, when $c=e=0$, it gives $S_{j,0}\le a+b$ and
$J_{j,0}\le a$.

One cannot simply assume $|T|+|G|\le15$. On a transitive set of fifteen
target rows, give every row the two child states $(0,2)$, where states
$0,1,2$ denote $(P_g,P_Q)=(0,0),(1,0),(1,1)$. Add a gate row $(2,2)$
pointing to all targets. Give all target edges label zero and all gate
edges label zero except its edge to the last target, labelled one.
There are fifteen target candidates and one gate candidate; their full
sixteen-set is incoherent at that last head. Likewise, three gate
candidates can form a directed cycle. These examples delimit the local
counting argument; they are not asserted to be full regular recursive
seeds or to attain all global profile maxima simultaneously.
\end{remark}
\section{Propagation, balancing, and a fixed-seed comparison}
\subsection{The inherited OR and composition principles}

\begin{proposition}[Seed and amplification]\label{prop:seed}
Suppose $s_0,s_1$ bound $S_{d,0},S_{d,1}$. For an integer $C\ge1$,
the OR of $C$ disjoint copies of $P_{d,1}$ is a total nonconstant
function $f$, rejects zero, and has
\begin{equation}\label{eq:seed}
 n=CLM(kr)^d,\qquad s(f)\le A=\max\{Cs_0,s_1\},\qquad
 \bs(f,0)\ge B=CMk^d.
\end{equation}
If $B>A^2>1$, the $m$-fold disjoint iterate $F_m$ has dimension $n^m$
and satisfies, for $\alpha=\log B/\log A>2$,
\begin{equation}\label{eq:amplify}
 F_m(0)=0,\quad s(F_m)\le A^m,\quad
 \bs(F_m,0)\ge B^m\ge s(F_m)^\alpha.
\end{equation}
In particular $\bs(F_m,0)$ tends to infinity.
\end{proposition}
\begin{proof}
At a zero of the OR every copy rejects, so sensitivities sum to at
most $Cs_0$. At a one, two accepting copies prevent any single-bit
rejection; with exactly one accepting copy, the sensitivity is at
most $s_1$. Lemma~\ref{lem:blocks} supplies $CMk^d$ disjoint
zero-sensitive blocks and proves nonconstancy.

For total $f,g$, a bit of child $i$ is sensitive for $f\circ g$
exactly when that bit is sensitive for $g$ and outer coordinate $i$
is sensitive for $f$ at the vector of child outputs. Hence
$s(f\circ g)\le s(f)s(g)$. If $g(0)=0$, disjoint zero-sensitive
families $U_1,\ldots,U_a$ for $f$ and $V_1,\ldots,V_b$ for $g$
give the nonempty, pairwise disjoint blocks $U_i\times V_j$ for
$f\circ g$. Flipping one changes exactly the outer coordinates in
$U_i$. Thus $\bs(f\circ g,0)\ge\bs(f,0)\bs(g,0)$.
Iteration proves~\eqref{eq:amplify}, and $B>1$ gives growth.
\end{proof}

\begin{proposition}[Published coarse consequence with explicit labels]
\label{prop:coarse}
For each $d\ge1$, choose
\begin{equation}\label{eq:qdepth}
 q(d)=\max\left\{3,1+\ceil{\frac{3^{d+1}}{13}}\right\},\qquad
 M=[13(q(d)-1)]^2.
\end{equation}
Then the explicit seed with $C=M$ satisfies
\begin{equation}\label{eq:coarsebounds}
 n=(d+2)M^2(kr)^d,\qquad
 A=2(d+2)M^{d+1},\qquad B=M^2k^d,
\end{equation}
and $B/A^2\ge2^d/[4(d+2)^2]$. At $d=9$ this is at least
$128/121>1$, so Proposition~\ref{prop:seed} gives a uniform power
separation. These quantitative estimates are inherited
from~\cite{OpenAI132}; explicit labels are supplied by
Theorem~\ref{thm:generator}.
\end{proposition}
\begin{proof}
Equation~\eqref{eq:qdepth} gives $M\ge9^{d+1}$. Normalize the
published profiles by $M^{j+b}$ and let $u_j,v_j$ be the maxima of
the normalized ordinary and joint profiles. Put
$\varepsilon=\max\{r,16\}/M\le2/3^{d+1}$.
Equations~\eqref{eq:ones} and~\eqref{eq:published} give
\[
 u_j\le(1+\varepsilon)u_{j-1}+3v_{j-1},\qquad
 v_j\le\varepsilon u_{j-1}+3v_{j-1},
\]
with $u_0\le L,v_0=0$. The corresponding equality comparison
sequences $U_j,V_j$ satisfy simultaneously
$U_j\le2L$ and $V_j\le\varepsilon L(3^j-1)$ for $j\le d$.
Indeed the bound for $V_j$ follows from the preceding bounds, and
summing the increments for $U$ gives
\[
 U_j\le L+\varepsilon L\left[\frac32(3^j-1)-j\right]<2L,
\]
since $\varepsilon\le2/3^{d+1}$. Thus
$S_{j,b}\le2LM^{j+b}$. Substituting in~\eqref{eq:seed} with $C=M$
gives~\eqref{eq:coarsebounds}, and $k/M^2>2$ gives the ratio.
\end{proof}

\subsection{The source-derived baseline used here}\label{sec:baseline}
Before comparing numerical consequences of Theorem~\ref{thm:profiles},
we include the following immediate improvements of the source proof
in the baseline. Use $S_{0,0}\le d+1$ rather than $d+2$, and
$|T|\le15$ rather than $16$, still charging every target repair to
$a$. One or two gate candidates cost at most $b+e$; three cost $3e$.
This gate envelope also appears in the zero-side recurrence
of~\cite{ProveItI}.
For simultaneous acceptance all gates cost $e$. Thus the immediate
source-derived zero-side updates are
\begin{equation}\label{eq:immediate}
 S_{j,0}\le15a+\max\{b+e,3e\},\qquad
 J_{j,0}\le15a+3e.
\end{equation}
Together with~\eqref{eq:ones}, these are valid for every $M\ge3$
with admissible labels. The large-$M$ restriction in
Proposition~\ref{prop:coarse} is only sufficient for its convenient
scalar estimate, not a restriction on exact recurrence iteration.
More generally, if the labels have coherence cap $K\ge3$, the identical
source-derived argument gives
\begin{equation}\label{eq:immediateK}
 S_{j,0}\le Ka+\max\{b+e,3e\},\qquad
 J_{j,0}\le Ka+3e.
\end{equation}
We use $K=7$ in both~\eqref{eq:immediateK} and
Corollary~\ref{cor:cap-profiles} for the fixed seed below.

The OR-copy count can also be optimized without a new combinatorial
theorem. If $W=Mk^d>s_0$, the function
\[
 \frac{\log(CW)}{\log\max\{Cs_0,s_1\}}
\]
increases up to $C=s_1/s_0$ and decreases afterward, where its
denominator is positive. Below that point the denominator is constant;
above it the derivative has the sign of $\log(s_0/W)<0$.
The two positive integers bracketing $s_1/s_0$ therefore suffice.
Both immediate refinements and OR optimization belong to the baseline.

\subsection{Certified comparison on one unchanged seed}\label{sec:numerical}
The following finite result compares proved upper bounds on sensitivity,
not exact sensitivity values. It changes neither the labels nor the
Boolean function between the two analyses.

\begin{theorem}[Fixed explicit seed]\label{thm:numerical}
Set
\begin{equation}\label{eq:fixed}
 \begin{gathered}
 q=1001,\quad d=17,\quad M=169000000,\quad r=13000,\\
 k=57122000000000001,\qquad C=164900179.
 \end{gathered}
\end{equation}
Use Theorem~\ref{thm:generator}'s labels and the recursion above.
The seed dimension and zero-block bound are exactly
\begin{equation}\label{eq:fixednb}
 n=C\cdot19M(kr)^{17},\qquad B=CMk^{17}.
\end{equation}
The integer $n$ has $373$ decimal digits. Its labels have exact coherence
cap seven by Corollary~\ref{cor:cyclic}. Iteration of the immediate
baseline~\eqref{eq:immediateK} with $K=7$ gives a bound $A_\old$;
iteration of Corollary~\ref{cor:cap-profiles} with the same $K=7$ gives
a strictly smaller bound $A_\new$.
For $\alpha_*=\log B/\log A_*$, certified enclosures are
\begin{align}
 2.015411246837677710&<\alpha_\old<2.015411246837677711,
 \label{eq:oldinterval}\\
 2.015411547511314020&<\alpha_\new<2.015411547511314021.
 \label{eq:newinterval}
\end{align}
Composition of this one explicit seed therefore satisfies
\eqref{eq:amplify} with $\alpha=\alpha_\new>2.0154$.
The strict improvement is from $A_\old$ to $A_\new$ at fixed $B,n$;
the inequality $\alpha>2.0154$ itself already holds for the baseline.
\end{theorem}

\begin{proof}[Proof and arithmetic certificate]
The exact definitions of the two integers are short recurrences.
Start with $v_0=(18,19M,0,0)$. For $v_j=(a,b,c,e)$, set
\begin{equation}\label{eq:update}
 v_{j+1}=\bigl(z,\ M^2a+rb,\ \min\{w,z\},\
              \min\{M^2c+rb,M^2a+rb\}\bigr),
\end{equation}
where $(z,w)$ is either the right side of~\eqref{eq:immediateK} with $K=7$
or $(F_0^{(7)},F_J^{(7)})$. Iterate exactly seventeen times. If $(a_*,b_*)$
are the first two final entries, define
$A_*=\max\{Ca_*,b_*\}$. At the terminal level the actual joint
profiles are zero; their final comparison entries are not used.
All operations are on integers and all cappings are valid by
Definition~\ref{def:profiles}.

Evaluation of these recurrences gives the exact rational enclosures
\begin{align*}
 3186883\cdot10^{143}&<A_\old<3186884\cdot10^{143},\\
 3186720\cdot10^{143}&<A_\new<3186721\cdot10^{143}.
\end{align*}
In particular $1<A_\new<A_\old$. For both values $A_*$, exact integer
powers verify
\begin{equation}\label{eq:powercert}
 A_*^{20154}<B^{10000}\le A_*^{20155}.
\end{equation}
Hence both exponent bounds exceed $2.0154$, and the common positive
$B$ with $A_\new<A_\old$ proves a strict improvement.

For the finer enclosures, the independent certificate uses exact rational
series bounds, without a library logarithm. For an integer $x>1$, put
$m=\lfloor\log_2x\rfloor$ and $z=(x-2^m)/(x+2^m)\in[0,1/3)$.
Then
\[
 \log x=m\log2+2\sum_{j=0}^{119}\frac{z^{2j+1}}{2j+1}+R,
 \qquad 0\le R\le\frac{2z^{241}}{241(1-z^2)}.
\]
The same series with $z=1/3$ encloses $\log2$. The implementation rounds
every operation outward on the rational lattice $10^{-100}\mathbb Z$.
Lower-numerator/upper-denominator and upper-numerator/lower-denominator
division then enclose each exponent ratio. The resulting intervals lie
strictly inside~\eqref{eq:oldinterval}--\eqref{eq:newinterval}.
The companion \nolinkurl{checks/independent-audit.json} records the exact
integers, longer intervals, and power checks~\eqref{eq:powercert}.
Equation~\eqref{eq:update} and parameters~\eqref{eq:fixed} independently
specify the certificate.
Finally Theorem~\ref{thm:generator} supplies the labels, and
Proposition~\ref{prop:seed} proves the stated Boolean-function bounds.
\end{proof}

No numerical search is needed to define this seed. The chosen $C$ was
selected by comparing the two bracketing integer candidates for the
earlier cap-fifteen immediate baseline. It is kept unchanged here so that
the Boolean function, block witnesses, and dimension remain identical.
It is not asserted to optimize either cap-seven bound. There is no claim
of a globally optimal exponent, minimal dimension, or simultaneous
attainability of all profile bounds. Larger exponents for the same
question are reported in~\cite{ProveItI,ProveItII}.

\begin{remark}[Comparison with the earlier certificate]
The first version used cap fifteen in both analyses. Deleting the seven
common leading zero digits already gives cap eight by
Theorem~\ref{thm:digits}; Corollary~\ref{cor:cyclic} improves this to the
exact cap seven. The following table applies each cap symmetrically, with
all seed parameters including $C$ fixed. Entries abbreviate certified
intervals; they are not exact sensitivities.
\begin{center}
\small
\begin{tabular}{@{}ccc@{}}
\toprule
Cap $K$ & Immediate baseline exponent & Coupled-profile exponent\\
\midrule
15 & $2.015410903091582180\ldots$ & $2.015411547311478476\ldots$\\
8 & $2.015411203868314436\ldots$ & $2.015411547486334582\ldots$\\
7 & $2.015411246837677710\ldots$ & $2.015411547511314020\ldots$\\
\bottomrule
\end{tabular}
\end{center}
At cap seven the gain is approximately $3.0067363631\times10^{-7}$.
The sharpened immediate baseline still falls below the earlier cap-fifteen
coupled bound, so the easy label tightening does not subsume the profile
refinement. The earlier \nolinkurl{checks/seed-certificate.json} is retained
and checked independently, including its original Decimal logarithm
enclosures~\cite{Decimal}.
\end{remark}

\section{Uniformity, verification, and limitations}
The seed is described by $(d,q,C)$ and a fixed program, requiring
$O(\log d+\log q+\log C)$ parameter bits. There is no label advice.
For~\eqref{eq:qdepth} with $C=M$, one can give only $d$ and derive the
other integers: $\log M=O(d)$ and $\log n=O(d^2)$. The parameter
description then has $O(\log d)$ bits, while deriving its large integers
takes time polynomial in $d$, not necessarily in $\log d$.

Polynomial logarithmic edge access is distinct from evaluation of the
whole $n$-bit function. A bottom-up evaluator computes every child's
predicate vector once. There are $O(n)$ nodes and at most
$d k(r+M^2)=O(dk^2)$ condition tests per internal node. Base distances
can be obtained by scanning each leaf once, since
\[
 D(x)=\wt(x)+L-2\max_v\wt(x|_{E_v}).
\]
A conservative bound is $O(ndk^2(\log n)^2)$ bit operations and
$O(n(d+1))$ bits for input, predicate values, and bookkeeping. Thus
expansion of an evaluation circuit has polynomial cost in the input
size; neither the description nor the generator requires a $2^n$-entry
truth table.

The companion files in \nolinkurl{topics/sensitivity-separation/} of the
\href{https://github.com/shjvandekamp/math}{repository} include the
generator, local profile diagnostics, exact recurrence implementation,
and the original and revised arithmetic certificates. The bounded generator check
exhausts $129088$ small tournaments and also checks a sharp coherent
fifteen-set for arbitrary orientations and parameter arithmetic. This is
compatible with the smaller cap in the indexed cyclic family.
The profile checker exhausts small
nested-state models and checks explicit candidate-count templates.
These are finite regression checks; the infinite label theorem and
all-input profile theorem rest on the proofs above. The fixed-seed
comparison is additionally an exact arithmetic calculation. A single
reproduction command from that topic directory is
\begin{center}
\texttt{python3 scripts/reproduce.py --fetch-sources}.
\end{center}
It verifies pinned source hashes and runs the bounded checks serially,
including a separate arithmetic audit of both the current cap-seven and
earlier cap-fifteen comparisons. That implementation
also checks the seven- and fourteen-vertex witnesses and uses the rational
logarithm enclosures above. It materializes no large tournament or Boolean
truth table. The independent arithmetic audit can also be run on its own:
\begin{center}
\texttt{python3 checks/independent\_audit.py}.
\end{center}

The scope of the contribution is limited. The first-difference technique,
the gated recursion, OR balancing, and composition are inherited methods.
The label theorem supplies an elementary directed adaptation with exact
uniform parameters. The profile theorem refines a local step of the
published proof. Part of its target/gate counting also appears
in~\cite{ProveItII}, as recorded in Section~1. No priority guarantee, record
exponent, or optimality claim is made.

\section*{AI assistance and verification status}
The arguments, code, and draft were developed with OpenAI Codex using
multiple GPT-6 Astra agents for construction, adversarial checking, and
writing. Separate agent reviews checked the saved mathematical arguments
and code, and the companion reproduction run passed. This describes
AI-assisted verification, not human peer review. The results of this note
have not been formalized in Lean, and no Lean kernel replay is claimed.
The written proofs and identified arithmetic artifacts state the evidence
supporting the conclusions.

\begin{thebibliography}{99}
\bibitem{AmbainisSun}
A. Ambainis and X. Sun,
\emph{New separation between $s(f)$ and $bs(f)$},
Electronic Colloquium on Computational Complexity, TR11-116 (2011),
\href{https://arxiv.org/abs/1108.3494v1}{arXiv:1108.3494v1};
see Section 4, Lemma 1 for the OR identities.

\bibitem{ColeVishkin}
R. Cole and U. Vishkin,
\emph{Deterministic coin tossing with applications to optimal parallel
list ranking}, Information and Control \textbf{70} (1986), no. 1, 32--53,
\href{https://doi.org/10.1016/S0019-9958(86)80023-7}{doi:10.1016/S0019-9958(86)80023-7}.

\bibitem{GMSW}
L. Gishboliner, A. Milojevi\'c, B. Sudakov, and Y. Wigderson,
\emph{Canonical Ramsey numbers of sparse graphs}, 11 October 2024,
\href{https://arxiv.org/abs/2410.08644v1}{arXiv:2410.08644v1};
see Section 2.1, the proof of Theorem 2.3.

\bibitem{Meiburg}
A. Meiburg,
\emph{Block Sensitivity can exceed Spectral Sensitivity Squared},
1 August 2026,
\href{https://arxiv.org/html/2608.00851v1}{arXiv:2608.00851v1}.

\bibitem{OpenAI132}
OpenAI,
\emph{A superquadratic separation between sensitivity and block
sensitivity}, 25 September 2026, OpenAI math family 132.
Pinned repository revision
\texttt{1581d4cbdd69a0faf34b2089614fe2a814316a6a}.
\href{https://github.com/openai/math/tree/1581d4cbdd69a0faf34b2089614fe2a814316a6a/preprints/A-superquadratic-separation-between-sensitivity-and-block-sensitivity-September-25-2026}{Manuscript and source files}.

\bibitem{ProveItI}
V. Reshetnikov (ProveIt repository; research companion prepared with ChatGPT),
\emph{An explicit exponent above $65/32$ for ordinary sensitivity versus
block sensitivity}, 7 October 2026, unrefereed. Placed at commit
\texttt{5476940edc3a2e7ea0c5426f013c6b5966451d73},
\href{https://github.com/VladimirReshetnikov/ProveIt/tree/5476940edc3a2e7ea0c5426f013c6b5966451d73/Combinatorics/BooleanFunctions/Research/sensitivity-block-sensitivity}{repository directory}.

\bibitem{ProveItII}
V. Reshetnikov (ProveIt repository; AI-assisted research draft),
\emph{Threshold-Resolved Amplification for Boolean Sensitivity}, draft dated
8 October 2026, unrefereed. Placed 9 October 2026 at commit
\texttt{18abb81b20977338c9c3c6511eaf333561709eca},
\href{https://github.com/VladimirReshetnikov/ProveIt/tree/18abb81b20977338c9c3c6511eaf333561709eca/Combinatorics/BooleanFunctions/Research/sensitivity-block-sensitivity}{repository directory}.

\bibitem{Decimal}
Python Software Foundation,
\emph{Decimal---Decimal fixed-point and floating-point arithmetic},
Python 3 documentation, consulted 15 October 2026.
\href{https://docs.python.org/3/library/decimal.html#decimal.Decimal.ln}{Correctly rounded logarithms}
and \href{https://docs.python.org/3/library/decimal.html#rounding-modes}{rounding modes}.

\bibitem{Tal}
A. Tal,
\emph{Properties and Applications of Boolean Function Composition},
Proceedings of ITCS 2013, pp. 441--454,
\href{https://doi.org/10.1145/2422436.2422485}{doi:10.1145/2422436.2422485}.

\bibitem{VishkinNotes}
U. Vishkin,
\emph{Thinking in Parallel: Some Basic Data-Parallel Algorithms and
Techniques}, 12 October 2010, p. 72,
\href{https://users.umiacs.umd.edu/~vishkin/PUBLICATIONS/classnotes.pdf}{author-hosted notes};
consulted 15 October 2026 for the position-plus-own-bit formulation.
\end{thebibliography}
\end{document}
